% Vietnamese translation for OI Public Library
% Source-File: IOI中国国家候选队论文/2009/骆可强/论程序底层优化的一些方法与技巧.ppt
\documentclass[11pt]{article}
\usepackage{oipl-vn}

\newcommand{\Oh}{\mathcal{O}}

\title{Bàn về một số phương pháp và kỹ thuật tối ưu tầng thấp của chương trình\\{\large Ghi chú theo slide}}
\author{Luo Keqiang\\Trường Trung học số 7 Thành Đô}
\date{Luận văn Trại đông 2009}

\begin{document}
\maketitle

\origtitle{Low-level optimization methods and techniques}
\sourcepath{IOI-China-Candidate-Team-Papers/2009/Luo-Keqiang/low-level-optimization-slides.ppt}

\section{Quan điểm}

Slide nhấn mạnh rằng thuật toán vẫn là yếu tố quyết định, nhưng khi độ phức
tạp đã phù hợp, chi tiết tầng thấp có thể tạo khác biệt lớn. Các kỹ thuật được
trình bày xoay quanh chi phí lệnh CPU, cache, dự đoán nhánh, phép toán số học
và biểu diễn bit.

\section{Các nhóm kỹ thuật}

\begin{itemize}
  \item Giảm phép chia hoặc thay bằng phép nhân, dịch bit khi điều kiện cho
  phép.
  \item Tổ chức dữ liệu theo thứ tự truy cập để tận dụng cache.
  \item Giảm rẽ nhánh khó dự đoán trong vòng lặp nóng.
  \item Dùng nén bit hoặc SIMD khi cùng một thao tác áp dụng cho nhiều phần
  tử nhỏ.
  \item Tránh chi phí định địa chỉ không cần thiết trong mảng nhiều chiều.
\end{itemize}

\section{Cách đọc ví dụ}

Phần ví dụ dẫn nhập trong bài viết đầy đủ cho thấy từng thay đổi nhỏ được đo
trên cùng một nền tảng. Khi đọc slide, nên quan sát quy trình: đo điểm nghẽn,
đưa ra giả thuyết về chi phí phần cứng, thay đổi cục bộ, rồi đo lại.

\section{Ghi chú}

Các kỹ thuật tầng thấp phụ thuộc máy và trình biên dịch. Vì vậy chúng nên được
dùng sau khi thuật toán đã đúng, có bộ kiểm thử ổn định, và có số đo chứng
minh phần được tối ưu thật sự nằm trên đường chạy nóng.

\end{document}
